\chapter{Синтез: что доказала серия \texttt{exp2--exp9a} и куда она вывела}

\section{Четыре слоя выводов}

Чтобы не смешивать разные типы результата, полезно собрать серию по слоям.

\subsection{Слой 1: инфраструктура и эмбрион}

Этот слой строят прежде всего \texttt{exp2} и \texttt{exp5}.

\begin{itemize}[leftmargin=2em]
\item Реальный эмбриональный граф можно собрать и прогонять через retrieval
pipeline.
\item Реальные источники порождают собственный класс ошибок: URI-labeling,
\texttt{RelatedTo}-hubness, embedding-lookup mismatch, ATOMIC fan-out.
\item Write-path требует batch-aware commit semantics; наивный single-pass
ломается даже на малом сценарии.
\end{itemize}

\subsection{Слой 2: локальная механика retrieval}

Этот слой задается \texttt{exp3}, \texttt{exp4}, \texttt{exp6}, \texttt{exp6b},
\texttt{exp6c}.

\begin{itemize}[leftmargin=2em]
\item Простой diffusion ranking действительно может предпочесть широкий хаб
специфическому ответу.
\item Semantic gate исправляет это лишь частично.
\item Причина не сводится к плохой similarity: topology of convergence сама по
себе умеет задавить более релевантный узел.
\item Post-accumulation correction имеет физический смысл.
\item Weighted leakage полезен, но не исчерпывает картину.
\item Capacity как ``эффективный объем представления'' выполняет реальную
работу на диагностических стендах.
\end{itemize}

\subsection{Слой 3: глобальная структурная слава}

Этот слой кристаллизуется в \texttt{exp6d}.

\begin{itemize}[leftmargin=2em]
\item Local response и global fame -- разные механизмы.
\item Мягкий PageRank prior действительно помогает, если применяется поверх
локально осмысленного response.
\item Полная PMI-подобная сила prior оказывается слишком агрессивной.
\end{itemize}

\subsection{Слой 4: перенос на реальный эмбрион}

Этот слой строят \texttt{7a--9a}.

\begin{itemize}[leftmargin=2em]
\item Automatic anchoring -- самостоятельный bottleneck.
\item Controlled-source regime нужен, но даже он не возвращает separation между
формулами автоматически.
\item Graph-native answer manifold на реальном субстрате часто не совпадает с
ручным canonical target.
\item Neither denominator tweaking, nor overlap, nor blind semantic-edge repair
дают решающего breakthrough.
\end{itemize}

\section{Что именно можно считать уже опровергнутым}

Серия достаточно уверенно отсекла несколько популярных, но ложных надежд.

\begin{enumerate}[leftmargin=2em]
\item \textbf{``Надо просто добавить semantic gate''} -- нет, этого мало.
\item \textbf{``Проблема только в embeddings''} -- нет, topology of convergence
делает самостоятельную работу.
\item \textbf{``Достаточно подобрать хороший denominator''} -- на реальном
эмбрионе это не разделяет методы так, как на золотых мирах.
\item \textbf{``Нужно просто поощрять overlap cue-fields''} -- нет, overlap не
создает путь там, где его нет.
\item \textbf{``Можно достроить все мосты blind kNN edges и получить решение''}
-- нет, тогда память начинает подтягивать морфологическое облако и опасно
съезжает к vector-RAG.
\end{enumerate}

\section{Что можно считать реально поддержанным}

Здесь нужна аккуратность. Поддержанность не означает окончательную истину, но
означает, что серия дала сильный положительный или диагностически устойчивый
сигнал в пользу гипотезы.

\begin{enumerate}[leftmargin=2em]
\item \textbf{Typed graph matters.} Тип ребра меняет физику retrieval;
\texttt{RelatedTo} нельзя считать нейтральным проводником.
\item \textbf{Post-accumulation viewpoint is correct.} Если дефект заключается в
слиянии множества вкладов на broad hub, штрафовать нужно именно merged
destination response.
\item \textbf{Local and global corrections are orthogonal.} Local retention и
PageRank-like prior решают разные задачи.
\item \textbf{Source injection must be modeled explicitly.} На реальном графе
anchoring не является технической мелочью; он архитектурно равноправен diffusion
и scoring.
\item \textbf{Target manifold is the right late-stage evaluation object.} Когда
графовая память возвращает соседний, но осмысленный узел, ошибочно считать это
тем же типом провала, что и победу пустого хаба.
\end{enumerate}

\section{Где сейчас настоящий bottleneck}

Если собрать все вместе, то bottleneck выглядит уже не как ``плохая формула''.
Скорее он состоит из трех связанных уровней.

\subsection{(1) Качество графового субстрата}

Эмбрион сейчас слишком смешан:

\begin{itemize}[leftmargin=2em]
\item lexical artifacts;
\item morphologically noisy neighborhoods;
\item неоднородные типы отношений;
\item неполные мосты между реально близкими представлениями;
\item слишком рыхлая интеграция episode shell.
\end{itemize}

\subsection{(2) Качество source injection}

Даже controlled-source regime понадобился не случайно. Это значит, что архитектура
естественного query-to-anchor mapping еще не дошла до состояния, где последующие
эксперименты можно интерпретировать без оговорок.

\subsection{(3) Предел самого spreading-подхода}

Есть все больше оснований думать, что для части задач память хочет не просто
``куда растечется активация'', а что-то ближе к:

\begin{itemize}[leftmargin=2em]
\item path-sensitive retrieval;
\item typed route preference;
\item consolidation shortcuts;
\item maybe small proof-like or script-like selection поверх поля активации.
\end{itemize}

\section{Почему серия все же продвинулась далеко}

Иногда длинная череда отрицательных и диагностических результатов производит
впечатление застоя. Но в этой серии произошло обратное: исследовательский объект
стал значительно яснее.

В начале речь шла почти обо всем сразу: реальный граф, embeddings, retrieval,
write-path, gate, priors. К концу серии можно уже сказать заметно точнее:

\begin{itemize}[leftmargin=2em]
\item проблема не в том, что ``ничего не работает'';
\item проблема не сводится к одному плохо выбранному коэффициенту;
\item есть рабочие локальные механизмы;
\item есть рабочие диагностические стенды;
\item есть честно исчерпанные ветки;
\item и есть новая, уже более содержательная развилка для следующего цикла.
\end{itemize}

\section{Что логично делать дальше}

Из самой серии вытекают по меньшей мере три разумные программы продолжения.

\subsection{Программа A: мета-аудит и смена класса retrieval}

Если принять, что spreading-подход почти уперся в потолок, следующий шаг --
формализовать это и перейти к path-selective retrieval. Тогда новая серия должна
прямо спрашивать:

\begin{quote}
какая причинно-смысловая траектория лучше всего объясняет запрос, а не какой
узел просто ярче всех после диффузии.
\end{quote}

\subsection{Программа B: архитектура консолидации эмбриона}

\texttt{9a} очень сильно намекнул, что графу не хватает внутренних мостов, но
достраивать их blind kNN-способом нельзя. Значит, нужна не просто edge repair, а
полноценная программа консолидации:

\begin{itemize}[leftmargin=2em]
\item alias-links;
\item synonym bridges;
\item shortcut links по часто используемым путям;
\item controlled reverse-links;
\item co-activation based strengthening;
\item фильтрация морфологического мусора.
\end{itemize}

\subsection{Программа C: развести слои памяти}

Еще один естественный ход -- перестать ожидать от одного плоского графа, что он
одновременно будет:

\begin{itemize}[leftmargin=2em]
\item словарем;
\item commonsense memory;
\item episode memory;
\item ассоциативной сетью синонимов и вариантов;
\item рабочей средой диффузии.
\end{itemize}

Тогда следующая архитектура должна быть многослойной, а transport и retrieval --
межслойными.

\section{Финальная оценка серии}

Если сформулировать общий итог одной фразой, он звучит так:

\begin{quote}
серия \texttt{exp2--exp9a} не нашла окончательную retrieval law для RIN, но она
очень качественно расчленила проблему и показала, где заканчивается сила
нынешнего диффузионного подхода и откуда должен начинаться следующий архитектурный
шаг.
\end{quote}

Это уже большой результат. Более того, без него любой следующий шаг рисковал бы
снова раствориться в хаотическом подборе коэффициентов и частных эвристик.

\begin{takeaway}
Самое ценное, что дала серия, -- не отдельная формула, а новая ясность. Теперь
можно продолжать не с надеждой, что ``еще один знаменатель все спасет'', а с
довольно четким пониманием: где нам нужен другой механизм чтения, где -- другой
механизм роста графа, а где -- более зрелая организация самой памяти как
многослойной архитектуры.
\end{takeaway}
